\subsection{Restricted Dual of an Algebra}

We assume from now on that K is a field.

For \(a \in A\) , denote by \(\boldsymbol{\gamma}(a)\) the mapping \(x \mapsto ax\) and by \(\boldsymbol{\delta}(a)\) the mapping \(x \mapsto xa\) from A to A. Then \(\boldsymbol{\gamma}\) is a linear representation of A in A, called the left regular representation; likewise, \(\delta\) is a linear representation of \(A^{\circ}\) in A, called (by abuse of language) the right regular representation of A. By transposition, we deduce from \(\delta\) and \(\gamma\) linear representations of the algebras A and \(A^{\circ}\) in the vector space \(A^{*}\) that define on \(A^{*}\) a left A-module structure and a left \(A^{\circ}\) -module structure. These two structures correspond to an (A, A)-bimodule structure on \(A^{*}\) with external laws defined by the formulas

\[
\langle a f, b \rangle = \langle f, b a \rangle ,\tag{3}
\]

\[
\langle f a, b \rangle = \langle f, a b \rangle\tag{4}
\]

for \(a, b, \in \mathrm{A}\) and \(f \in \mathrm{A}^*\) .

Recall (II, §2, No.~4, p.~234, Definition 4) that if E is a linear subspace of A, then its orthogonal \(E'\) in \(A^{*}\) is the set of linear forms on A whose restriction to E is zero. Likewise, the orthogonal \(F'\) in A of a linear subspace F of \(A^{*}\) is the intersection of the kernels of the linear forms belonging to F.

We know (II, §7, No.~5, p.~299, Theorem 7) that the mapping \(\varphi : \mathrm{E} \mapsto \mathrm{E}'\) is a bijection from the set of linear subspaces of A of finite codimension to the set of linear subspaces of \(\mathrm{A}^*\) of finite dimension. The inverse mapping \(\varphi^{-1}\) sends a finite-dimensional subspace F of \(\mathrm{A}^*\) to its orthogonal \(\mathrm{F}'\) in A.

PROPOSITION 1. --- a) The mapping \(\varphi\) induces a bijection from the set of left (resp. right, two-sided) ideals of A of finite codimension to the set of right A-submodules (resp. left A-submodules, sub-bimodules) of \(A^{*}\) of finite dimension over K.

\begin{enumerate}
\def\labelenumi{\alph{enumi})}
\setcounter{enumi}{1}
\tightlist
\item
  Every left or right ideal of A of finite codimension contains a two-sided ideal of finite codimension.
\end{enumerate}

Formulas (3) and (4) prove that if E is a left ideal of A, then \(E'\) is a right A-submodule of \(A^{*}\) ; if V is a right A-submodule of \(A^{*}\) , then \(V'\) is a left ideal of A. The cases of right ideals and of two-sided ideals are treated analogously. This proves a).

Let \(\mathfrak{a}\) be a left ideal of A of finite codimension. The left A-module \(\mathrm{E} = \mathrm{A}_s / \mathfrak{a}\) is finite-dimensional over K, and so is \(\mathrm{End}_{\mathrm{K}}(\mathrm{E})\) . The kernel of the homomorphism \(a \mapsto a_{\mathrm{E}}\) from A to \(\mathrm{End}_{\mathrm{K}}(\mathrm{E})\) is therefore a two-sided ideal of A of finite codimension contained in \(\mathfrak{a}\) . Passing to the opposite ring \(\mathrm{A}^{\circ}\) , we see that every right ideal of A of finite codimension contains a two-sided ideal of finite codimension.

DEFINITION 2. --- The restricted dual

of the K-algebra A, denoted by \(\Theta(A)\) , is the union in \(A^{*}\) of the orthogonals of the two-sided ideals of \(A\) of finite codimension.

By Proposition 1, we can give the following equivalent descriptions of \(\Theta(A)\) :

- the union of the orthogonals of the left ideals of A of finite codimension

-- the union of the orthogonals of the right ideals of A of finite codimension

- the union of the orthogonals of the two-sided ideals of \(\mathrm{A}\) of finite codimension

- the union of the left A-submodules of \(\mathrm{A}^*\) of finite dimension over \(\mathbf{K}\)

- the union of the right A-submodules of \(\mathrm{A}^*\) of finite dimension over \(K\)

- the union of the \((\mathrm{A},\mathrm{A})\) -sub-bimodules of \(\mathrm{A}^*\) of finite dimension over \(\mathbf{K}\) .

We have \(\Theta(A) = \Theta(A^{\circ})\) , and \(\Theta(A) = A^{*}\) if \(A\) has finite degree over \(K\) . Let \(f \in A^{*}\) ; then \(f\) belongs to \(\Theta(A)\) if and only if \(Af\) (resp. \(fA\) , \(AfA\) ) is a linear subspace of \(A^{*}\) of finite dimension over \(K\) .

The sum of two (A, A)-sub-bimodules of \(A^{*}\) of finite dimension over K has finite dimension over K. It follows that \(\Theta(\mathrm{A})\) is an (A, A)-sub-bimodule of \(A^{*}\) .

